\section{Closed elements and topological completeness}
\label{sec:p08-mt-topological-completeness}

McKinsey and Tarski's \emph{On Closed Elements in Closure Algebras}
(1946, pp.~122--162) develops a closed-element algebra, not a second
definition of the modern Heyting implication. Its Theorems~1.14--1.19
connect that algebra with closure and topology; Theorem~1.11 and
Lemmas~2.22--2.24 provide the finite bridge; Theorem~3.25 gives the
functional freeness of Euclidean closed-set algebras
\cite{McKinseyTarski1946}. We first reconstruct the closed-element
construction, then give explicit finite spaces and a real-line realization.
The latter realization is a separate Cantor-interval reconstruction of the
Euclidean conclusion, not a transcription of the paper's dissectability
argument or of its 1944 predecessor.

\paragraph{The direction of the algebra matters.}
The Heyting algebra of Definition~\ref{def:p08-heyting-algebra} orders truth
values upwards towards \(1\). The source's closed-element algebra is its
order dual. We retain the source's lattice operations \(+\) and \(\cdot\).
We render its binary Brouwerian difference as \(a\mathbin{\dot{-}}b\),
distinguishing it from Boolean complementation \(-b\); this is a typed
transcription of the source operation, not the modern implication.

\begin{hbdefinition}{def:p08-mt-brouwerian}{The 1946 closed-element convention}
A Brouwerian algebra here is a bounded distributive lattice
\(A=(L,+,\cdot,0,1)\) with an operation \(\dot{-}\) satisfying
\[
 a\mathbin{\dot{-}}b\leq c
 \quad\Longleftrightarrow\quad a\leq b+c.
\]
Its Brouwerian complement is \(\mathord{\sim}a=1\mathbin{\dot{-}}a\).
This is the convention of Definition~1.1 and Theorem~1.3 of the source
(pp.~124--125), nowadays also called a bounded co-Heyting algebra.
For a closure algebra, the Boolean complement is instead written \(-a\).
\end{hbdefinition}

\begin{hbdefinition}{def:p08-mt-source-functions}{Uniform Brouwerian-algebraic functions}
In the source's Definition~2.2 (pp.~136--137), an \(n\)-argument
Brouwerian-algebraic function is a family of operations, one on each
Brouwerian algebra, generated by the \(n\) projections and the constant
unit function using \(+\), \(\cdot\), and Brouwerian difference.
Equality of such functions means equality for all assignments in all
these algebras, not just equality of two functions on one selected space.
The zero function is expressible as \(1\mathbin{\dot{-}}1\).
Equivalently, each finite expression in those operations supplies such
a uniform family. This algebraic presentation is distinct from a
sequent calculus or a chosen list of IPL axiom schemata.
\end{hbdefinition}

\begin{hblemma}{thm:p08-mt-dual}{proved=annex:p08-mt.1}{The typed Heyting--Brouwerian bridge}
If \(H\) is a Heyting algebra, its order dual is a Brouwerian algebra with
\[
 a+_A b=a\wedge_H b,\quad a\cdot_A b=a\vee_H b,
 \quad0_A=1_H,\quad1_A=0_H,
 \quad a\mathbin{\dot{-}}_A b=b\Rightarrow_H a.
\]
Conversely these equations reconstruct \(H\) from \(A\).
In particular a formula has Heyting value \(1_H\) exactly when its
closed-convention term has value \(0_A\), not \(1_A\).
\end{hblemma}

\begin{hbdefinition}{def:p08-mt-closure}{Closure algebra and its closed elements}
A closure algebra is a Boolean algebra \(B\) with a unary operation \(C\)
such that \(C0=0\), \(x\leq Cx\), \(CCx=Cx\), and
\(C(x+y)=Cx+Cy\). Write \(Ix=-C(-x)\).
Its closed elements satisfy \(Cx=x\). On them set
\[
 a\mathbin{\dot{-}}b=C(a\cdot(-b)).
\]
For a topological space, \(B=\mathcal P(X)\), \(C\) is closure, and
\(+\), \(\cdot\), \(-\) are union, intersection, and relative complement.
Neither a Hausdorff nor a separation axiom is imposed.
\end{hbdefinition}

\begin{hblemma}{thm:p08-mt-closed}{proved=annex:p08-mt.2}{Closed elements really have the required residual}
The closed elements of any closure algebra form a Brouwerian algebra.
For closed \(a,b,c\),
\[
 C(a\cdot(-b))\leq c\quad\Longleftrightarrow\quad a\leq b+c.
\]
For closed subsets of a space the difference is
\(\operatorname{cl}(a\setminus b)\). Complementation identifies their
order dual with the open-set Heyting algebra, whose implication is
\(U\Rightarrow V=\operatorname{int}((X\setminus U)\cup V)\).
This reconstructs Theorem~1.14 and the adjoining open-element observation
on p.~130 of the source.
\end{hblemma}

\paragraph{The Boolean extension is an actual construction.}
Suppose \(e:L\hookrightarrow B\) is a bounded lattice embedding into a
Boolean algebra. Let \(K\) be its Boolean-generated subalgebra. Every
element has a finite normal form
\[
 x=\sum_{i=1}^{m} e(a_i)\cdot(-e(b_i));
 \qquad
 Cx=\sum_{i=1}^{m} e(a_i\mathbin{\dot{-}}b_i).
\]
The empty sum represents zero. The formula for \(C\) is not merely a
name for a desired extension: independence from the normal form has to
be proved. This is the construction in Theorem~1.15 (pp.~130--132).

\begin{hbtheorem}{thm:p08-mt-extension}{proved=annex:p08-mt.3}{The minimal closed-element extension}
The displayed prescription is well defined. It makes \(K\) a closure
algebra whose closed elements are precisely \(e[L]\), with their original
Brouwerian difference. For every \(x\in K\), \(Cx\) is the least element
of \(e[L]\) above \(x\). Any isomorphism of Brouwerian algebras extends
uniquely to an isomorphism of these Boolean-generated closure extensions.
For arbitrary algebras, existence of the starting Boolean embedding can
be obtained from bounded distributive lattice separation; this use of
the Boolean prime ideal theorem is distinct from the construction above.
For the finite algebras used below, an explicit embedding needs no choice.
\end{hbtheorem}

\begin{hblemma}{thm:p08-mt-finite}{proved=annex:p08-mt.4}{Finite preservation of the values actually used}
Let \(M\) be a finite set of \(n\) elements of a Brouwerian algebra.
The bounded distributive sublattice \(D\) generated by \(M\) has at most
\(2^{2^n}\) elements. With its own least-difference operation it is a
Brouwerian algebra. Whenever \(a,b,a\mathbin{\dot{-}}b\) all belong to
\(D\), the new and old differences agree.
Consequently every finite list of terms, including all its subterms, can
be evaluated in a finite algebra with exactly the original values of
those terms. The generated \emph{lattice} is not asserted to be a
subalgebra for every original difference.
This expands the preservation clause of Theorem~1.11 (pp.~127--128) and
the finite-approximation argument of Lemma~2.22 (p.~144).
\end{hblemma}

\begin{hbdefinition}{def:p08-mt-connected}{Relativization and well-connectedness}
For \(a\in L\), put \(A_a=\{x:a\leq x\}\), with the old lattice
operations, zero \(a\), unit \(1\), and difference
\(x\mathbin{\dot{-}}_a y=a+(x\mathbin{\dot{-}}y)\).
A Brouwerian algebra is well-connected when
\(x\cdot y=0\) implies \(x=0\) or \(y=0\).
These are Definitions~1.6 and~1.20 of the source.
\end{hbdefinition}

\begin{hblemma}{thm:p08-mt-minimal}{proved=annex:p08-mt.5}{A minimal finite separating algebra is well-connected}
The map \(x\mapsto a+x\) is a Brouwerian homomorphism \(A\to A_a\).
If a term is nonzero somewhere, there is a finite well-connected algebra
and an assignment where it is nonzero. In particular, a formula which
fails in a Heyting algebra fails in the order dual of such an algebra.
This gives the reduction of Lemmas~2.23--2.24 (pp.~144--145), with the
relativization maps and the smaller-algebra argument made explicit.
\end{hblemma}

\begin{hbdefinition}{def:p08-mt-finite-space}{The finite closed-element space}
For a finite distributive lattice \(L\), let \(J(L)\) be its nonzero
join-irreducible elements, ordered as in \(L\), and put
\(e(a)=\{j\in J(L):j\leq a\}\).
Give \(J(L)\) the upper Alexandrov topology: its opens are all upsets.
Thus its closed sets are downsets and
\(\operatorname{cl}(S)=\mathord{\downarrow}S\).
\end{hbdefinition}

\begin{hbtheorem}{thm:p08-mt-finite-space}{proved=annex:p08-mt.6}{Representation with a separation witness}
The map \(e\) identifies \(L\) with all downsets of \(J(L)\), preserves
the Brouwerian difference, and generates the Boolean algebra
\(\mathcal P(J(L))\). Its closure is exactly the normal-form closure of
Theorem~\ref{thm:p08-mt-extension}. If \(a\not\leq b\), some
\(j\leq a\) satisfies \(j\not\leq b\), so the embedding separates
\(a\) from \(b\) at a specified point.
If \(L\) is nontrivial and well-connected, \(J(L)\) has a least point.
The one-element lattice has the empty representing space instead.
\end{hbtheorem}

\begin{hbdefinition}{def:p08-mt-formula-term}{Formula translation into the closed convention}
For the IPL grammar \(\bot,\wedge,\vee,\to\), let
\[
 t_p=a_p,\quad t_\bot=1,\quad
 t_{B\wedge D}=t_B+t_D,\quad t_{B\vee D}=t_B\cdot t_D,
 \quad t_{B\to D}=t_D\mathbin{\dot{-}}t_B.
\]
Negation abbreviates \(B\to\bot\). In the finite space set
\(\nu(p)=J(L)\setminus e(a_p)\).
\end{hbdefinition}

This translation connects the host's IPL language and proved Lindenbaum
algebra with the source's algebraic functions. It does not attribute the
host's LJ presentation, or a later prime-filter syntax, to the 1946 paper.

\begin{hblemma}{thm:p08-mt-formulas}{proved=annex:p08-mt.7}{Every connective, in both conventions}
For every formula \(B\),
\[
 \nu(B)=J(L)\setminus e(t_B).
\]
Thus \(t_B=0\) iff \(\nu(B)=J(L)\), and a witness
\(j\in e(t_B)\) is a point where \(B\) fails.
Equivalently, on any upper Alexandrov space,
\(w\in\nu(B\to D)\) iff for every \(v\geq w\),
\(v\in\nu(B)\) implies \(v\in\nu(D)\).
The finite algebra and topological interpretation preserve the formula,
not just its atomic values.
\end{hblemma}

\paragraph{A real space, rather than a change of vocabulary.}
To realize a finite rooted space on \(\mathbb R\), use the standard
middle-thirds Cantor set \(K\subseteq[0,1]\). Its complementary bounded
intervals have generations \(1,2,\ldots\). With \(k\) colors, assign
an interval of generation \(n\) the color \(1+((n-1)\bmod k)\).
The colors will correspond to the immediate successors of a root.

\begin{hblemma}{thm:p08-mt-cantor}{proved=annex:p08-mt.8}{Every color occurs arbitrarily near each Cantor point}
For \(k\geq1\), every open neighborhood of every \(x\in K\) contains
a whole complementary interval of each color. In particular it is not
enough for intervals merely to accumulate somewhere in \(K\): each
color has this property at every point, including \(0\) and \(1\).
\end{hblemma}

\begin{hbtheorem}{thm:p08-mt-realization}{proved=annex:p08-mt.9}{An explicit interior surjection onto each finite rooted space}
Let \(W\) be a finite nonempty poset with a least point \(r\), with the
upper Alexandrov topology. There is a continuous, open, surjective map
\(f:\mathbb R\to W\). For \(W=\{r\}\), take the constant map.
Otherwise map \(K\) to \(r\), color its complementary bounded
intervals by the finitely many immediate successors \(s\) of \(r\),
and on an interval of color \(s\) use a homeomorphism from that interval
to \(\mathbb R\), followed by the recursively constructed map onto
\(\uparrow s\). Map the two exterior intervals by the construction
for any one successor. The recursion terminates because
\(|\uparrow s|<|W|\).
\end{hbtheorem}

\begin{hblemma}{thm:p08-mt-pullback}{proved=annex:p08-mt.10}{Preservation under an interior map}
If \(f:X\to Y\) is continuous and open, then for every subset \(S\)
of \(Y\),
\[
 f^{-1}[\operatorname{int}_Y(S)]
 =\operatorname{int}_X(f^{-1}[S]),\qquad
 f^{-1}[\operatorname{cl}_Y(S)]
 =\operatorname{cl}_X(f^{-1}[S]).
\]
Pulling back an open valuation therefore preserves the value of every
IPL formula; pulling back closed values preserves every Brouwerian
term. If \(f\) is also surjective, both kinds of validity are reflected,
and a point of failure has a point of failure in its inverse image.
\end{hblemma}

\begin{hbtheorem}{thm:p08-mt-topological-completeness}{proved=annex:p08-mt.11}{The topological completeness promised in P8}
For every IPL formula \(B\), the following are equivalent:
\[
 \begin{array}{ll}
 \textnormal{(i)}&\vdash_{\mathrm{IPL}}B;\\
 \textnormal{(ii)}&B\textnormal{ is valid on all topological spaces};\\
 \textnormal{(iii)}&B\textnormal{ is valid on all finite upper Alexandrov spaces};\\
 \textnormal{(iv)}&B\textnormal{ is valid on }\mathbb R;\\
 \textnormal{(v)}&B\textnormal{ is valid on }\mathbb R^d
                    \textnormal{ for each fixed positive integer }d.
 \end{array}
\]
Here validity quantifies over all assignments of open sets to the
variables and means that the formula has value the entire space.
If \(B\) is not derivable, the proof constructs a finite algebra, a
nonempty rooted finite space, and an open valuation on \(\mathbb R\)
where every point of the Cantor set is a failure point. For
\(\mathbb R^d\), projection to the first coordinate transfers this
countervaluation. No finiteness of the underlying Euclidean space,
Hausdorff assumption on the finite space, or Boolean prime ideal theorem
is used in this finite-support argument.
\end{hbtheorem}

The algebraic input to this proof is the already proved IPL Lindenbaum
countervaluation in Theorem~\ref{thm:p10-heyting-completeness}.
The topological part is not left inside that algebraic theorem: finite
approximation, the closure construction, the pointwise witnesses, and
the real-line interior map are supplied above. This is theoremhood
completeness. It does not assert strong completeness of a fixed
Euclidean space for arbitrary infinite premise sets, or replace the
separate Rasiowa--Sikorski dense-set lemma.

\begin{hbcorollary}{thm:p08-mt-functional-free}{proved=annex:p08-mt.12}{The reconstructed Euclidean functional-freeness conclusion}
For every positive integer \(d\), the Brouwerian algebra of all closed
subsets of \(\mathbb R^d\) has exactly the identities common to all
Brouwerian algebras. Explicitly, if an equation \(s=t\) between finite
Brouwerian terms fails in some algebra, there is an assignment of closed
sets in \(\mathbb R^d\) for which it fails.
Thus this algebra is functionally free in the sense of source
Definition~3.19; this does not say that it is freely generated by a
finite chosen set, or that the Boolean extension of a free algebra is
functionally free as a closure algebra.
\end{hbcorollary}

The source proves a broader result: Theorem~3.25 also covers nonempty normal,
dense-in-itself spaces with a countable basis, via its dissectability
and universality results. The explicit Cantor construction above proves
the Euclidean instance and is not a proof of that entire broader class.
The source's extension of dissectable Brouwerian algebras to closure
algebras in Lemma~3.8 uses countably many copies and a saturated closure;
it is a different operation from replacing a finite order by its
Alexandrov topology. Its geometrical prerequisites from the 1944 paper
remain distinct from the local arguments here.

\begin{hblemma}{thm:p08-mt-boundaries}{proved=annex:p08-mt.13}{Non-Boolean examples and limits of the hypotheses}
On the two-point space \(W=\{r<s\}\), with opens
\(\varnothing,\{s\},W\), take \(\nu(p)=\{s\}\). Then
\(\nu(\neg p)=\varnothing\) and
\(\nu(p\vee\neg p)=\{s\}\ne W\).
On \(\mathbb R\), take \(\nu(p)=\mathbb R\setminus\{0\}\).
Then \(\nu(\neg p)=\varnothing\) and
\(\nu(\neg\neg p)=\mathbb R\), while excluded middle still fails at
\(0\). By contrast every discrete space has a Boolean open-set algebra
and validates excluded middle. The zero-dimensional Euclidean space
\(\mathbb R^0\), a singleton, is therefore not an IPL-characterizing
space, and cannot be included in the positive-dimensional conclusion.
\end{hblemma}

In the closed convention on the same two-point space, put \(a=\{r\}\).
Its Brouwerian complement is \(1\mathbin{\dot{-}}a
=\operatorname{cl}(\{s\})=W\), rather than the Boolean complement
\(\{s\}\). Hence \(a\cdot(1\mathbin{\dot{-}}a)=\{r\}\ne0\).
This displays the failure of the Boolean complement law in the native
closed-element algebra, alongside its open-set IPL interpretation.

\paragraph{Why the map hypotheses cannot be dropped.}
The identity from a discrete two-point space to the preceding
Sierpi\'nski space is a continuous surjection but is not open.
For the subset \(\{r\}\), its inverse image has nonempty interior,
whereas the inverse image of its target interior is empty. Thus mere
continuity cannot justify implication preservation. Conversely an
open map need not take inverse images of opens to opens; continuity
is also required. Surjectivity is what reflects a separating witness;
without it the image may miss every failure point.
